% Part: first-order-logic
% Chapter: introduction
% Section: models-theories

\documentclass[../../../include/open-logic-section]{subfiles}

\begin{document}

\olfileid[fr]{fol}{int}{mod}

\olsection{Modèles et théories}

Une fois la syntaxe et la sémantique de la logique du premier ordre
définies, nous pouvons étudier les propriétés des !!{structure}s et des
notions sémantiques. Nous pouvons aussi définir des systèmes de
!!{derivation} et les étudier. Étant donné un ensemble
d'!!{sentence}s, nous pouvons demander quelles !!{structure}s rendent
vrais tous les !!{sentence}s de cet ensemble. Si $\Gamma$ est un
ensemble d'!!{sentence}s, toute !!{structure}~$\Struct{M}$ qui les
satisfait est appelée un \emph{modèle de~$\Gamma$}. Nous pouvons partir
de~$\Gamma$ et chercher ses modèles : quelle forme ont-ils ? Quelle
taille minimale ou maximale peuvent-ils avoir ? Nous pouvons aussi
partir d'une seule !!{structure}, ou d'une collection de
!!{structure}s, et demander quels !!{sentence}s y sont vrais.
Existe-t-il des !!{sentence}s qui \emph{caractérisent} ces
!!{structure}s, au sens où ils sont vrais dans ces structures, et dans
elles seules ? Ces questions relèvent de la \emph{théorie des
modèles}. Elles sous-tendent également la \emph{méthode axiomatique},
qui décrit une collection de !!{structure}s au moyen d'un ensemble
d'!!{sentence}s, les axiomes d'une théorie. Cette méthode repose sur
l'observation que les !!{sentence}s entraînés en logique du premier
ordre par les axiomes sont exactement ceux qui sont vrais dans tous les
modèles de ces axiomes.

Prenons comme exemple très simple les préordres. Un préordre est une
relation~$R$ sur un ensemble~$A$ qui est à la fois réflexive et
transitive. Un ensemble~$A$ muni d'une relation binaire
$R \subseteq A \times A$ fournit exactement les données nécessaires à
!!a{structure} d'un langage du premier ordre ne possédant qu'un seul
symbole de relation binaire~$\Obj P$ : nous posons
$\Domain{M} = A$ et $\Assign{\Obj P}{M} = R$. Puisque $R$ est un
préordre, elle est réflexive et transitive, et nous pouvons trouver un
ensemble~$\Gamma$ d'!!{sentence}s de la logique du premier ordre qui
expriment ces deux propriétés :
\begin{align*}
  & \lforall[\Obj v_0][\Atom{\Obj P}{\Obj v_0,\Obj v_0}]\\
  & \lforall[\Obj v_0][\lforall[\Obj v_1][\lforall[\Obj v_2]
    [((\Atom{\Obj P}{\Obj v_0,\Obj v_1} \land
       \Atom{\Obj P}{\Obj v_1,\Obj v_2})
      \lif \Atom{\Obj P}{\Obj v_0,\Obj v_2})]]]
\end{align*}
Ces !!{sentence}s symbolisent respectivement « pour tout~$x$, $Rxx$ »
($R$ est réflexive) et « si $Rxy$ et $Ryz$, alors $Rxz$ » ($R$ est
transitive). Nous voyons qu'!!a{structure}~$\Struct{M}$ est un modèle
de ces deux !!{sentence}s~$\Gamma$ si et seulement si $R$, c'est-à-dire
$\Assign{\Obj P}{M}$, est un préordre sur~$A$, c'est-à-dire sur
$\Domain{M}$. Autrement dit, les modèles de~$\Gamma$ sont exactement
les préordres. Toute propriété commune à tous les préordres qui peut
s'exprimer dans le langage du premier ordre dont~$\Obj P$ est le seul
!!{predicate}, comme la réflexivité et la transitivité ci-dessus, est
entraînée par les deux !!{sentence}s de~$\Gamma$, et réciproquement.
Tout résultat démontré au sujet des modèles de~$\Gamma$ vaut donc pour
tous les préordres.

Pour chaque théorie et chaque classe de modèles particulières, comme
$\Gamma$ et la classe de tous les préordres, des questions intéressantes
se posent sur ce qui peut et ne peut pas s'exprimer dans le langage du
premier ordre correspondant. Certaines propriétés des !!{structure}s
présentent un intérêt pour tous les langages et toutes les classes de
modèles, notamment celles qui concernent la taille du !!{domain}. Pour
tout~$n \in \PosInt$, on peut toujours exprimer que le !!{domain}
contient exactement $n$~!!{element}s. On peut aussi exprimer, au moyen
d'un ensemble infini d'!!{sentence}s, que le !!{domain} est infini. En
revanche, on ne peut exprimer ni que le domaine est fini, ni qu'il est
!!{nonenumerable}. Ces limitations des langages du premier ordre
découlent des théorèmes de compacité et de L\"owenheim--Skolem.

\end{document}
